% 第22章《光的干涉》23道原书编号章末习题，独立推导候选解析。
% V2主PDF262--264，S辅导PDF312--318。原图、专家物理审核仍待办理。
\section*{第22章\quad 光的干涉：独立解析}

\paragraph{习题 22.1}
波列长度 \(L=c\tau\)，周期数 \(N=L/\lambda\)，故
\[
\boxed{N=\frac{c\tau}{\lambda}
=\frac{(2.998\times10^8)(10^{-8})}{589.3\times10^{-9}}
\simeq5.09\times10^6}.
\]
这里的波数指波列中完整周期个数，而非角波数 \(2\pi/\lambda\)。

\paragraph{习题 22.2}
傍轴双缝的相邻明纹间距
\(\Delta x=D\lambda/d\)，故
\[
\boxed{\lambda=\frac{d\Delta x}{D}
=\frac{(.60\times10^{-3})(2.27\times10^{-3})}{2.5}
=5.448\times10^{-7}\,\mathrm m\simeq545\,\mathrm{nm}}.
\]
这是绿色光，与汞灯546.1\,nm绿色谱线吻合到测量精度。

\paragraph{习题 22.3}
镜面虚像源 \(S'\) 位于镜下2\,mm，与实源相隔
\(d=2h=4\,\mathrm{mm}\)。屏到源投影距离
\(D=20+30=50\,\mathrm{cm}\)，不是仅镜长30\,cm。
镜面反射多出相对相位 \(\pi\)，所以镜面延长线
为\emph{暗纹}；第一明纹的垂直位移是普通明纹间距的一半：
\[
\boxed{x_1=\frac{D\lambda}{2d}
=\frac{(.50)(7.2\times10^{-7})}{2(.004)}
=4.5\times10^{-5}\,\mathrm m=45\,\mathrm{\mu m}}.
\]
图22.33的镜面有限长度会限制条纹可见范围，
本推导取傍轴虚光源近似。

\paragraph{习题 22.4}
第1到第10级明纹之间有 \(10-1=9\) 个间距：
\(\Delta x=36\,\mathrm{mm}/9=4.0\,\mathrm{mm}\)，
\[
\boxed{\lambda=\frac{d\Delta x}D
=\frac{(.15\times10^{-3})(4.0\times10^{-3})}{1}
=6.0\times10^{-7}\,\mathrm m=600\,\mathrm{nm}}.
\]

\paragraph{习题 22.5}
令从东方向向北偏的夹角为 \(\theta\)，
两发射台间距 \(d=3000\,\mathrm m\)，波长
\(\lambda=c/f\simeq1500\,\mathrm m\)。
按图中南台相位落后北台 \(\pi/2\)，
远场的相位差（取原书方位约定）
\[
\Delta\phi=2\pi\frac{d\sin\theta}{\lambda}+\frac\pi2.
\]
相长干涉要求 \(2\pi m\)，故
\[
\sin\theta=\frac{\lambda}{d}(m-\tfrac14)
=\tfrac12m-\tfrac18.
\]
由 \(|\sin\theta|\le1\)，整数
\(m=-1,0,1,2\)，得到
\[
\boxed{\theta=-38.7^\circ,\ -7.18^\circ,\
22.0^\circ,\ 61.0^\circ\ \text{(约)}}.
\]
这是向东半平面的四支；向西半平面对应镜像射线。
反转初始相位差的符号会反转有向角，故必须记载台站相位先后。

\paragraph{习题 22.6}
相邻静区中心对应两源路程差改变一波长。
在傍轴近似下 \(\lambda=d\Delta x/D\)：
\[
\lambda=\frac{12\times4.5}{50}\,\mathrm{cm}=1.08\,\mathrm{cm},
\qquad
\boxed{f=\frac{u}{\lambda}
=\frac{25}{1.08}\,\mathrm{Hz}
\simeq23.1\,\mathrm{Hz}}.
\]

\paragraph{习题 22.7}
入射斜光先给两缝引入 \(d\sin\varphi\)
的初始光程差；向 \(\theta\) 方向出射再引入
\(d\sin\theta\)。相对净光程差
\(\Delta=d(\sin\theta-\sin\varphi)\)，
因此
\[
\boxed{d\sin\theta-d\sin\varphi=m\lambda,\ m\in\mathbb Z}.
\]
小角度 \(\sin\theta\simeq\theta\) 时，
\(\theta_m\simeq\sin\varphi+m\lambda/d\)，
故
\[
\boxed{\Delta\theta=\theta_{m+1}-\theta_m\simeq\lambda/d},
\]
与 \(\varphi\) 无关。严格局部微分近似应为
\(\lambda/(d\cos\theta)\)，仅小 \(\theta\) 可忽略该因子。

\paragraph{习题 22.8}
直接波与海面镜面反射波的几何光程差为
\(\Delta l\simeq2h\sin\theta\)，
理想海面反射产生额外相位 \(\pi\)。
相消时
\[
\frac{2\pi}{\lambda}(2h\sin\theta)+\pi
=(2m+1)\pi
\quad\Longrightarrow\quad
2h\sin\theta=m\lambda.
\]
排除 \(m=0\) 的零掠射角后，最小正解取 \(m=1\)：
\[
\boxed{\theta_{\min}=\arcsin\frac{c}{2hf}
=\arcsin\frac{2.998\times10^8}{2(25)(6.0\times10^7)}
\simeq5.74^\circ}.
\]
海水非理想反射系数依偏振和入射角变化；
该结果限原图所用镜面模型。

\paragraph{习题 22.9}
两缝等振幅时
\(I(\theta)=I_{\max}\cos^2(\delta/2)\)，
\(\delta=2\pi d\sin\theta/\lambda\)。
半峰值处 \(\cos^2(\delta/2)=1/2\)，
在中心两侧 \(\delta=\pm\pi/2\)，
即 \(\theta_{\pm}\simeq\pm\lambda/(4d)\)。
于是两侧\emph{半峰值之间的全宽}
\[
\boxed{\Delta\theta_{\rm FWHM}
=\theta_+-\theta_-\simeq\frac{\lambda}{2d}}.
\]
若把单侧偏移称半宽，则应为 \(\lambda/(4d)\)；
本题按原书给出的两个半峰值角间隔解释。

\paragraph{习题 22.10}
图22.35两叠加平面波的相邻等相位线
在屏上的投影距离 \(\Delta x\)
满足 \(\Delta x\cos45^\circ=\lambda\)：
\[
\boxed{\Delta x=\frac{632.8\,\mathrm{nm}}{\cos45^\circ}
\simeq895\,\mathrm{nm}}.
\]
叠加振幅 \(A_1:A_2=2:1\)，亮暗峰值强度
分别 \(\propto(A_1+A_2)^2\)、\(\propto(A_1-A_2)^2\)。
因此
\[
\boxed{V=\frac{I_{\max}-I_{\min}}{I_{\max}+I_{\min}}
=\frac{2A_1A_2}{A_1^2+A_2^2}
=\frac45=0.80}.
\]
图示是特殊倾角平面波，不可误用普通双缝纹距公式。

\paragraph{习题 22.11}
用谱宽反比的数量级定义
\[
\boxed{\tau_c\simeq(\Delta\nu)^{-1}
=\frac1{1.3\times10^9}
\simeq7.69\times10^{-10}\,\mathrm s},
\quad
\boxed{L_c\simeq c\tau_c\simeq0.231\,\mathrm m}.
\]
具体线型与谱宽定义可影响系数；此处与教材取法一致，
不代表所有激光器精确相干函数均满足等号。

\paragraph{习题 22.12}
太阳作为角宽约 \(\alpha=10^{-2}\,\mathrm{rad}\)
的扩展非相干光源，各方向入射角跨过 \(\alpha\)
引起双缝相差量级 \(d\alpha\)。
干涉仍有可见度需 \(d\alpha\lesssim\lambda\)：
\[
\boxed{d_{\max}\sim\frac{\lambda}{\alpha}
=\frac{550\times10^{-9}}{10^{-2}}
=5.5\times10^{-5}\,\mathrm m=55\,\mathrm{\mu m}}.
\]
这是空间相干宽度估算，不是严格绝对截止阈值。

\paragraph{习题 22.13}
玻璃替代原来等长的空气，在一条缝增加
\((n-1)h\) 光程。第7级原亮纹移至原零级，
光程差对应七个波长：
\[
(n-1)h=7\lambda,\qquad
\boxed{h=\frac{7(550\,\mathrm{nm})}{1.58-1}
\simeq6.64\,\mathrm{\mu m}}.
\]

\paragraph{习题 22.14}
空气--SiO\(_2\) 与SiO\(_2\)--Si
两界面均为从低折射率向高折射率反射，
两个反射场\emph{均}有 \(\pi\) 相位反转；
相对反射相位差为零。
反射暗纹要求
\[
2nh=(k+\tfrac12)\lambda,\quad k=0,1,2,\ldots
\]
零厚楔尖为亮纹；从此起第7条暗纹 \(k=6\)，
故
\[
\boxed{h=\frac{13\lambda}{4n}
=\frac{13(589.3\,\mathrm{nm})}{4(1.50)}
\simeq1.28\,\mathrm{\mu m}}.
\]
因上、下反射均有半波损失，
不能误套空气--玻璃--空气薄膜的暗纹序号。

\paragraph{习题 22.15}
空气--玻璃--空气上界面反射有半波损失，
下界面反射无半波损失。
正入射反射加强
\(2nh=(k+\frac12)\lambda_{\rm R}\)，
透射加强 \(2nh=m\lambda_{\rm T}\)。
由 \(2nh=2(1.50)(0.4)=1.2\,\mathrm{\mu m}\)，
\[
\lambda_{\rm R}=\frac{4nh}{2k+1}
=\frac{2400\,\mathrm{nm}}{2k+1}
\Longrightarrow\boxed{\lambda_{\rm R}=480\,\mathrm{nm}\ (k=2)}.
\]
\[
\lambda_{\rm T}=\frac{2nh}{m}
=\frac{1200\,\mathrm{nm}}m
\Longrightarrow
\boxed{\lambda_{\rm T}=600\,\mathrm{nm}\ (m=2),\
400\,\mathrm{nm}\ (m=3)}.
\]
采用 \(400\,\mathrm{nm}\) 为可见下限时将其纳入；
更严格视觉阈值可能归为近紫外边界。忽略色散。

\paragraph{习题 22.16}
空气--SiO薄膜反射相位变 \(\pi\)，
SiO--水晶反射无相位反转；
反射增强最小膜厚是四分之一波厚：
\[
\boxed{h_{\min}=\frac{\lambda}{4n_f}
=\frac{560\,\mathrm{nm}}{4(2.0)}
=70\,\mathrm{nm}}.
\]
真实多层膜的反射率还受界面振幅比影响。

\paragraph{习题 22.17}
空气--油膜--玻璃的两次反射均在低到高
折射率界面，\emph{相对}反射相变为零。
反射消光级次满足
\[
2n_oh=(k+\tfrac12)\lambda.
\]
相邻两个消光波长 \(\lambda_1=485\,\mathrm{nm}\)
和 \(\lambda_2=679\,\mathrm{nm}\) 满足
\((k+\frac12)\lambda_1=(k-\frac12)\lambda_2\)。
两式消元：
\[
\boxed{h=\frac{\lambda_1\lambda_2}
{2n_o(\lambda_2-\lambda_1)}
=\frac{(485)(679)}
{2(1.32)(679-485)}\,\mathrm{nm}
\simeq643\,\mathrm{nm}}.
\]
“再次相消”必须指相邻消光级次；
否则存在整数倍厚度。

\paragraph{习题 22.18}
皂膜两侧均为空气，仅在第一次反射有半波
相位变化。入射角 \(45^\circ\)，
由Snell定律
\(n\sin r=\sin45^\circ\)。
第一个反射亮纹满足
\(2nh\cos r=\lambda/2\)，因此
\[
\boxed{h_{\min}=\frac{500\,\mathrm{nm}}
{4\sqrt{1.33^2-\sin^245^\circ}}
\simeq111\,\mathrm{nm}}.
\]
法线方向同一膜的增强波长
\[
\boxed{\lambda_\perp=4nh
\simeq4(1.33)(111\,\mathrm{nm})
\simeq590\,\mathrm{nm}},
\]
故约为黄色，而非原斜视方向的绿色。

\paragraph{习题 22.19}
空气--MgF\(_2\)--玻璃两界面反射
均有相位变化 \(\pi\)，故使两反射光
传播相位差为 \(\pi\) 即得到减反条件：
\[
\boxed{h_{\min}=\frac{\lambda}{4n_2}
=\frac{550\,\mathrm{nm}}{4(1.38)}
\simeq99.6\,\mathrm{nm}}.
\]
本膜折射率 \(1.38\) 与理想零反射幅值匹配
\(\sqrt{1.52}\simeq1.23\) 不同，
因此此厚度表示尽量相消，并不保证反射率严格为零。

\paragraph{习题 22.20}
牛顿环相邻明环的平方半径差为 \(R\lambda\)，
相隔五个级次有
\(r_{m+5}^2-r_m^2=5R\lambda\)。
直径 \(d=2r\)，得
\[
\boxed{\lambda=\frac{d_{m+5}^2-d_m^2}{20R}
=\frac{(4.60\times10^{-3})^2-(3.00\times10^{-3})^2}
{20(1.03)}
\simeq5.90\times10^{-7}\,\mathrm m
=590\,\mathrm{nm}}.
\]
平方差可消去未知起始级次，故不需判定第几环。

\paragraph{习题 22.21}
单程玻璃替换空气增加光程
\((n-1)h\)。迈克耳孙仪同一臂往返各穿一次，
所以两臂光程差变化为
\[
\boxed{\Delta(\mathrm{OPD})=2(n-1)h}.
\]
若两臂同时插相同薄片，
公共增加量不产生相对光程差。

\paragraph{习题 22.22}
平面反射镜移动 \(\Delta L\) 使往返光程差
改变 \(2\Delta L\)。每过一条中心干涉纹改变一个波长，
\[
N\lambda=2\Delta L,\quad
\boxed{\lambda=\frac{2(0.3220\times10^{-3})}{1204}
=5.349\times10^{-7}\,\mathrm m
=534.9\,\mathrm{nm}}.
\]

\paragraph{习题 22.23}
按原书图22.37，两路径分别\emph{单程}穿过两个等长
气室；充入一室的气体取代真空引入
\(\Delta(\mathrm{OPD})=(n-1)d\) 的路差。
当 \(N=98\) 条干涉纹移动，
\[
N\lambda=(n-1)d,\qquad
\boxed{n=1+\frac{N\lambda}{d}
=1+\frac{98(589.3\times10^{-9})}{.20}
\simeq1.000289}.
\]
若误用22.21那种同一臂往返通过玻璃的系数2，
会把气体的折射率增量少算一半。
